The concentration sensitivity of Au, Ni, Cu, Pd, Pt, and Ti was examined at 298.15~K by varying one concentration at a time. The nominal aqueous metal-ion concentration was varied from $10^{-6}$ to $10^{-2}\,\mathrm{M}$, the total ammonia concentration from $0.002$ to $0.1\,\mathrm{M}$, and the total glycine concentration from $0.01$ to $1.0\,\mathrm{M}$. Other concentrations were held at their reference values: nominal aqueous metal-ion concentration $10^{-4}\,\mathrm{M}$, $[\ce{NH3}]_{\mathrm{tot}}=0.02\,\mathrm{M}$, and $[\mathrm{Gly}]_{\mathrm{tot}}=0.1\,\mathrm{M}$, without added cyanide. The nominal metal-ion concentration corresponds to the activity input under the ideal-solution convention. Each black outline marks pH 10.5--13.5 and $-2.0$ to $2.3$~V vs RHE, converted to SHE and clipped to the plotted range.

% Au sensitivity panels
\begin{figure}[!htbp]
    \centering
    \includegraphics[width=\textwidth]{Au_activity_three_levels_1x3_side_colors.pdf}
    \caption{Au Pourbaix diagrams with the nominal aqueous metal-ion concentration varied from $10^{-6}$ to $10^{-2}\,\mathrm{M}$. Other concentrations were held at their reference values.}
    \label{fig:Au_activity_sensitivity_SI}
\end{figure}

\begin{figure}[!htbp]
    \centering
    \includegraphics[width=\textwidth]{Au_NH3_three_levels_1x3_side_colors.pdf}
    \caption{Au Pourbaix diagrams with the total ammonia concentration varied from $0.002$ to $0.1\,\mathrm{M}$. Other concentrations were held at their reference values.}
    \label{fig:Au_NH3_sensitivity_SI}
\end{figure}

\begin{figure}[!htbp]
    \centering
    \includegraphics[width=\textwidth]{Au_Gly_three_levels_1x3_side_colors.pdf}
    \caption{Au Pourbaix diagrams with the total glycine concentration varied from $0.01$ to $1.0\,\mathrm{M}$. Other concentrations were held at their reference values.}
    \label{fig:Au_Gly_sensitivity_SI}
\end{figure}



% Ni sensitivity panels
\begin{figure}[!htbp]
    \centering
    \includegraphics[width=\textwidth]{Ni_activity_three_levels_1x3_side_colors.pdf}
    \caption{Ni Pourbaix diagrams with the nominal aqueous metal-ion concentration varied from $10^{-6}$ to $10^{-2}\,\mathrm{M}$. Other concentrations were held at their reference values.}
    \label{fig:Ni_activity_sensitivity_SI}
\end{figure}

\begin{figure}[!htbp]
    \centering
    \includegraphics[width=\textwidth]{Ni_NH3_three_levels_1x3_side_colors.pdf}
    \caption{Ni Pourbaix diagrams with the total ammonia concentration varied from $0.002$ to $0.1\,\mathrm{M}$. Other concentrations were held at their reference values.}
    \label{fig:Ni_NH3_sensitivity_SI}
\end{figure}

\begin{figure}[!htbp]
    \centering
    \includegraphics[width=\textwidth]{Ni_Gly_three_levels_1x3_side_colors.pdf}
    \caption{Ni Pourbaix diagrams with the total glycine concentration varied from $0.01$ to $1.0\,\mathrm{M}$. Other concentrations were held at their reference values.}
    \label{fig:Ni_Gly_sensitivity_SI}
\end{figure}



% Cu sensitivity panels
\begin{figure}[!htbp]
    \centering
    \includegraphics[width=\textwidth]{Cu_activity_three_levels_1x3_side_colors.pdf}
    \caption{Cu Pourbaix diagrams with the nominal aqueous metal-ion concentration varied from $10^{-6}$ to $10^{-2}\,\mathrm{M}$. Other concentrations were held at their reference values.}
    \label{fig:Cu_activity_sensitivity_SI}
\end{figure}

\begin{figure}[!htbp]
    \centering
    \includegraphics[width=\textwidth]{Cu_NH3_three_levels_1x3_side_colors.pdf}
    \caption{Cu Pourbaix diagrams with the total ammonia concentration varied from $0.002$ to $0.1\,\mathrm{M}$. Other concentrations were held at their reference values.}
    \label{fig:Cu_NH3_sensitivity_SI}
\end{figure}

\begin{figure}[!htbp]
    \centering
    \includegraphics[width=\textwidth]{Cu_Gly_three_levels_1x3_side_colors.pdf}
    \caption{Cu Pourbaix diagrams with the total glycine concentration varied from $0.01$ to $1.0\,\mathrm{M}$. Other concentrations were held at their reference values.}
    \label{fig:Cu_Gly_sensitivity_SI}
\end{figure}



% Pd sensitivity panels
\begin{figure}[!htbp]
    \centering
    \includegraphics[width=\textwidth]{Pd_activity_three_levels_1x3_side_colors.pdf}
    \caption{Pd Pourbaix diagrams with the nominal aqueous metal-ion concentration varied from $10^{-6}$ to $10^{-2}\,\mathrm{M}$. Other concentrations were held at their reference values.}
    \label{fig:Pd_activity_sensitivity_SI}
\end{figure}

\begin{figure}[!htbp]
    \centering
    \includegraphics[width=\textwidth]{Pd_NH3_three_levels_1x3_side_colors.pdf}
    \caption{Pd Pourbaix diagrams with the total ammonia concentration varied from $0.002$ to $0.1\,\mathrm{M}$. Other concentrations were held at their reference values.}
    \label{fig:Pd_NH3_sensitivity_SI}
\end{figure}

\begin{figure}[!htbp]
    \centering
    \includegraphics[width=\textwidth]{Pd_Gly_three_levels_1x3_side_colors.pdf}
    \caption{Pd Pourbaix diagrams with the total glycine concentration varied from $0.01$ to $1.0\,\mathrm{M}$. Other concentrations were held at their reference values.}
    \label{fig:Pd_Gly_sensitivity_SI}
\end{figure}



% Pt sensitivity panels
\begin{figure}[!htbp]
    \centering
    \includegraphics[width=\textwidth]{Pt_activity_three_levels_1x3_side_colors.pdf}
    \caption{Pt Pourbaix diagrams with the nominal aqueous metal-ion concentration varied from $10^{-6}$ to $10^{-2}\,\mathrm{M}$. Other concentrations were held at their reference values.}
    \label{fig:Pt_activity_sensitivity_SI}
\end{figure}

\begin{figure}[!htbp]
    \centering
    \includegraphics[width=\textwidth]{Pt_NH3_three_levels_1x3_side_colors.pdf}
    \caption{Pt Pourbaix diagrams with the total ammonia concentration varied from $0.002$ to $0.1\,\mathrm{M}$. Other concentrations were held at their reference values.}
    \label{fig:Pt_NH3_sensitivity_SI}
\end{figure}

\begin{figure}[!htbp]
    \centering
    \includegraphics[width=\textwidth]{Pt_Gly_three_levels_1x3_side_colors.pdf}
    \caption{Pt Pourbaix diagrams with the total glycine concentration varied from $0.01$ to $1.0\,\mathrm{M}$. Other concentrations were held at their reference values.}
    \label{fig:Pt_Gly_sensitivity_SI}
\end{figure}



% Ti sensitivity panels
\begin{figure}[!htbp]
    \centering
    \includegraphics[width=\textwidth]{Ti_activity_three_levels_1x3_side_colors.pdf}
    \caption{Ti Pourbaix diagrams with the nominal aqueous metal-ion concentration varied from $10^{-6}$ to $10^{-2}\,\mathrm{M}$. Other concentrations were held at their reference values.}
    \label{fig:Ti_activity_sensitivity_SI}
\end{figure}

\begin{figure}[!htbp]
    \centering
    \includegraphics[width=\textwidth]{Ti_NH3_three_levels_1x3_side_colors.pdf}
    \caption{Ti Pourbaix diagrams with the total ammonia concentration varied from $0.002$ to $0.1\,\mathrm{M}$. Other concentrations were held at their reference values.}
    \label{fig:Ti_NH3_sensitivity_SI}
\end{figure}

\begin{figure}[!htbp]
    \centering
    \includegraphics[width=\textwidth]{Ti_Gly_three_levels_1x3_side_colors.pdf}
    \caption{Ti Pourbaix diagrams with the total glycine concentration varied from $0.01$ to $1.0\,\mathrm{M}$. Other concentrations were held at their reference values.}
    \label{fig:Ti_Gly_sensitivity_SI}
\end{figure}



